\section{Implementation and Code Snippets}
\label{sec:impl}

All models, tasks and kernels were implemented for this project in PyTorch and Triton
\citep{tillet2019triton}; no external state-space or linear-attention library is used. The code is available
at \url{https://github.com/Devx228/EE798-project}. The main files are \texttt{layers/ssm.py} (selective and
rotational SSM), \texttt{layers/deltanet.py} (Gated DeltaNet and DeltaProduct), \texttt{layers/attention.py},
\texttt{layers/triton\_scan.py}, \texttt{tasks/} (data generators), \texttt{train.py} and the experiment
scripts. The code was written with the help of an AI coding assistant; the experiments were run on the
author's own GPU, and all reported numbers come from those runs.

\paragraph{Reference recurrence.} Listing~\ref{lst:delta} is the token-by-token form of
\eqref{eq:delta}. It is slow, but it is the ground truth for testing the chunked algorithm.
\begin{lstlisting}[caption={Reference gated delta rule (\texttt{layers/deltanet.py}).},label={lst:delta}]
for t in range(T):
    S = S * torch.exp(g[:, :, t])[..., None, None]          # decay: alpha_t * S
    v_old = torch.einsum("bhkv,bhk->bhv", S, k[:, :, t])    # what memory returns for k_t
    S = S + torch.einsum("bhk,bhv->bhkv", k[:, :, t],
                         beta[:, :, t, None] * (v[:, :, t] - v_old))  # erase + write
    outs.append(torch.einsum("bhkv,bhk->bhv", S, q[:, :, t]))         # o_t = S^T q_t
\end{lstlisting}
The first line applies the decay $\alpha_t=e^{g_t}$ to the whole state. The second reads the value currently
stored under the key, $S^\top k_t$. The third adds $\beta_t k_t(v_t-S^\top k_t)^\top$, which is the
gradient step of Section~\ref{sec:math}: with $\beta_t=1$ the old value is fully replaced. The last line
reads the memory with the query.

\paragraph{Changing the eigenvalue range.} The negative-eigenvalue variant is a one-line change after the
discretisation \eqref{eq:disc} (Listing~\ref{lst:neg}); for DeltaNet the analogous change is
\texttt{beta = 2 * beta}.
\begin{lstlisting}[caption={Decay of the selective SSM (\texttt{layers/ssm.py}).},label={lst:neg}]
dt = F.softplus(dt_raw + self.dt_bias)        # step size Delta_t > 0
a = torch.exp(-dt * torch.exp(self.A_log))    # a_t = exp(Delta_t * A), A < 0, so a_t in (0, 1)
if self.neg_eigen:
    a = 2.0 * a - 1.0                         # a_t in (-1, 1): the state can flip sign
\end{lstlisting}

\paragraph{Rotations as a data-dependent RoPE.} Equation \eqref{eq:rope} is implemented literally: the
cumulative angle is a prefix sum, $B$ and $C$ are rotated by it, and the verified real scan does the rest.
\begin{lstlisting}[caption={Rotational SSM (\texttt{layers/ssm.py}).},label={lst:rot}]
def rotational_scan(x, a, theta, B, C, chunk_size=64):
    phase = theta.cumsum(dim=1)                       # Phi_t = sum_{k<=t} theta_k
    return ssd_scan_chunked(x, a, rotate_pairs(B, -phase), rotate_pairs(C, -phase), chunk_size)
\end{lstlisting}

\paragraph{DeltaProduct by interleaving.} Rather than writing a new kernel, the $n_h$ Householder steps of
\eqref{eq:product} are laid out along the time axis: token $t$, sub-step $j$ becomes position $tn_h+j$.
The decay is applied only at the first sub-step and the query only at the last, so the verified chunked
delta rule computes DeltaProduct exactly.
\begin{lstlisting}[caption={DeltaProduct via interleaving (\texttt{layers/deltanet.py}).},label={lst:prod}]
q_x = q.new_zeros(b, T, n_h, H, d); q_x[:, :, -1] = q   # read memory after the last sub-step
g_x = g.new_zeros(b, T, n_h, H);    g_x[:, :, 0] = g    # decay once per token
o, S = delta_rule_chunked(*flatten(q_x, k, v, beta, g_x))  # sequence of length T * n_h
o = o[:, :, n_h - 1 :: n_h]                               # one output per token
\end{lstlisting}

\paragraph{Triton scan kernel.} Listing~\ref{lst:triton} is the inner loop of the forward kernel. One GPU
program handles one (sequence, head, block of $P$ channels) and keeps its $[B_P\times N]$ slice of the state
in registers for the whole sequence, so the state never travels to global memory; only $x,a,B,C$ are read
and $y$ is written.
\begin{lstlisting}[caption={Inner loop of the Triton scan kernel (\texttt{layers/triton\_scan.py}).},label={lst:triton}]
state = tl.zeros([BLOCK_P, BLOCK_N], dtype=tl.float32)
for t in range(T):
    a_t = tl.load(a_base + t * s_at)
    x_t = tl.load(x_base + t * s_xt, mask=mask_p, other=0.0)
    b_t = tl.load(b_base + t * s_bt, mask=mask_n, other=0.0)
    c_t = tl.load(c_base + t * s_ct, mask=mask_n, other=0.0)
    state = a_t * state + x_t[:, None] * b_t[None, :]     # h_t = a_t h_{t-1} + x_t B_t^T
    y_t = tl.sum(state * c_t[None, :], axis=1)             # y_t = h_t C_t
    tl.store(y_base + t * s_yt, y_t, mask=mask_p)
\end{lstlisting}

\paragraph{Data generation.} The MQAR generator writes $n$ key--value pairs at the start of the sequence and
places each key once more as a query at a random later position, followed by its value; only query positions
are scored. Keys are drawn without replacement from the lower half of the vocabulary and values from the
upper half. The group-word generator samples uniform group elements and computes prefix products with a
precomputed multiplication table (for $S_3$, the composition table of all six permutations).

\paragraph{Verification.} A test suite (62 tests on the CPU plus the Triton kernel tests on the GPU) checks that (i) every chunked algorithm (SSM with positive
and negative decays, rotational SSM, gated delta rule with $\beta\in(0,2)$, DeltaProduct) matches its
token-by-token reference in float64 (maximum error $\approx10^{-15}$) for sequence lengths that are and are not
multiples of the chunk length; (ii) the single-token decoding of every model reproduces its parallel forward
pass; (iii) the Triton kernel matches the reference (run on the GPU, and on CPU in Triton's interpreter mode);
and (iv) the task generators produce correct labels and the group tables satisfy the group axioms. Two
further tests check the algebra behind the extensions: a rotation by $2\pi/3$ returns the state after three
steps, and the product of two reflections is a rotation of order three when the keys are $60^\circ$ apart.
